\documentclass[instructions.tex]{subfiles}
\begin{document}
 
\chapter{Il faut aimer ses ennemis et leur pardonner.}
 
Jésus-Christ nous commande, sous peine de damnation, d’aimer nos ennemis, et de leur pardonner. Nous ne pouvons nous plaindre de ce commandement, si nous regardons trois choses : ce que nous sommes, qui est notre ennemi, et le mal qu’on nous fait.
 
\primo Nous sommes pécheurs ; nous avons besoin de la miséricorde de Dieu. Sans cette miséricorde, toutes les créatures s’armeroient pour nous détruire, comme des ennemis de leur créateur\footnote{Misericordiæ Domini, quia non sumus consumpti. Thr. 3.}. Etant pécheurs, nous méritons d’être réprouvés de Dieu. Toute la différence qu’il y a entre nous et un réprouvé, c’est qu’un réprouvé n’a ni pardon ni miséricorde à espérer, au lieu que nous espérons le pardon et la miséricorde de Dieu. Ne serions-nous donc pas bien malheureux de refuser le pardon à nos ennemis ? Il est juste que Dieu nous traite comme nous aurons traité nos frères. Si nous refusons de pardonner, nous nous fermons le ciel, et nous devons nous attendre à être jugés sans miséricorde\footnote{Judicium enim sine misericordiâ illi qui non fecit misericordiam. Jac. 2.}.
 
\secundo Qui est notre ennemi ? c’est notre prochain comme notre frère ; c’est-à-dire un homme, un enfant de Dieu, comme nous, qui est aimé de Dieu, peut-être plus que nous, que Dieu prend sous sa sauve-garde, et auquel il nous défend d’attenter. Notre prochain est comme le lieutenant de Jésus-Christ. Tout ce que vous lui faites, dit le Sauveur, fût-il le dernier des hommes, je le tiens fait à moi-même\footnote{Quamdiù fecistis uni ex his fratribus meis minimis, mihi fecistis. Matth. 25.}. C’est donc Jésus-Christ que vous haïssez, c’est donc contre Jésus-Christ que vous exercez votre fureur, lorsque vous haïssez votre prochain. Si votre ennemi mérite votre haine, Jésus-Christ ne la mérite pas ; c’est néanmoins sur le Sauveur que retombe votre haine, lorsque vous haïssez vos frères.
 
\tertio Mais avez-vous bien examiné pourquoi vous haïssez quelqu’un ? c’est peut-être parce qu’il vous a averti de vos défauts ou des désordres de votre famille, parce qu’il vous a empêché d’avoir un parti ou un emploi dont vous êtes indigne. Si cela est, il vous a rendu service ; vous devriez l’en estimer ; si vous en concevez de la haine, c’est un effet de votre mauvais fond et de votre malice.
 
De quoi êtes-vous offensé ? d’un rien, d’une parole échappée, d’un rapport qui est faux et que vous avez cru ; et voilà un effet de votre légèreté ou de votre mauvais génie. Votre ennemi n’est peut-être coupable que dans vos idées. Pendant la nuit, dit le prophète, les bêtes sauvages et les monstres sortent, mais au lever du soleil ils disparaissent. Dans votre humeur noire, votre esprit est comme dans une nuit ; il se remplit de nuages et d’idées monstrueuses ; il s’effarouche, se défie et s’ombrage de tout. Prévenu contre une personne, vous vous imaginez qu’elle a pensé, qu’elle a dit, qu’elle a fait, ce qu’elle n’a jamais dit ni jamais fait. Vous croyez qu’elle cherche à se faire valoir, à vous supplanter, à vous faire peine, quoiqu’elle n’y pense pas ; c’est que les ténèbres sont dans votre esprit. Mais laissez entrer le soleil de la vérité et la charité dans votre cœur ; expliquez-vous avec cette personne, parlez-lui avec ouverture de cœur ; et bientôt toutes vos pensées noires, et ces monstres qui vous déchirent disparaitront.
 
Quand même votre prochain vous auroit offensé et traité aussi cruellement qu’on a traité Jésus-Christ, devriez-vous pour cela le haïr ? Jésus-Christ n’a-t-il pas embrassé Judas qui le trahissoit, n’a-t-il pas offert le pardon à ses bourreaux ? Seriez-vous assez dur pour résister à l’exemple d’un Dieu, qui pardonne sa mort ? Vous voulez vous venger ? Ah ! chrétien, regardez et écoutez votre Sauveur attaché à la croix\footnote{Vis vindicari ? Vide pendentem, audi precantem. S. Aug.}. Votre prochain vous a-t-il offensé ? mais n’avez-vous jamais offensé Dieu ni fait de peine à personne ? et quand vous n’auriez jamais offensé le prochain, n’avez-vous jamais offensé Dieu ni mérité l’enfer ? ô homme de peu de foi ! comprenez que, quand un ennemi auroit enlevé vos biens, attenté à votre vie, noirci votre réputation, tout cela n’est rien en comparaison de ce que vous devez à la justice de Dieu ; cependant, quelque immenses que soient vos dettes envers Dieu, si vous pardonnez, Jésus-Christ vous promet que tout vous sera pardonné\footnote{Dimitte et dimittentur. Matth.}. Celui, dit saint Augustin, qui refuse d’acquitter ses péchés par un moyen si facile et si court, n’a plus rien à espérer pour le ciel.
 
Saint Jean Gualbert, chevalier, ayant pardonné à un mortel ennemi, dont il étoit facile de se venger, alla tout de suite faire à Dieu cette prière : Seigneur, vous avez promis le pardon à celui qui pardonne ; vous savez, ô mon Dieu, les péchés dont je suis coupable : je viens vous supplier et vous sommer de tenir votre parole et de me pardonner, puisque je viens de pardonner à mon ennemi, pour l’amour de vous. Il sentit dans le moment, avec consolation, que sa prière étoit exaucée. Il vécut dans la sainteté, et fonda un ordre célèbre.
 
Celui qui se venge, fait une action basse et indigne, parce qu’il se rend esclave de sa passion. Mais celui qui pardonne, fait une action pleine de gloire, parce qu’il triomphe de lui-même.
 
\section{Résolutions}
 
\primo Quand j’aurai sujet de me plaindre de mon prochain, je me rappellerai que Jésus-Christ m’ordonne de l’aimer, qu’il m’en a donné le premier l’exemple ; que les saints l’ont imité, et que je me ferme à jamais l’entrée du ciel, si je n’accomplis pas sincerement ce précepte.
 
\secundo En conséquence, je pardonnerai généreusement les offenses qu’on m’aura faites. Mon Dieu, remplissez mon cœur de miséricorde envers mon prochain, afin que je puisse réciter, sans prononcer ma propre condamnation, cette prière que vous nous avez prescrite : Pardonnez-nous nos offenses, comme nous les pardonnons à ceux qui nous ont offensés.
 
\end{document}
